\section{OpenAI 收入账：看得见的收入与看不见的潜在收入}

上一章解释了为什么巨头不容易下牌桌，本章把镜头收回到 OpenAI。围绕 OpenAI 最大的争议，是巨额算力承诺与未来现金流是否匹配。广密的拆法很有教学意义：先分“看得见的收入”，再分“看不见的收入”。前者像互联网平台和云服务，后者则取决于 Agent 能否替代劳动、AI 能否做出人类做不了的高价值任务。

\begin{figure}[H]
\centering
\includegraphics[width=0.92\textwidth]{openai-revenue-stack.png}
\caption{OpenAI 收入栈：订阅、广告电商、API、Agent 劳动与 AI 科学家。自制概念图，依据 00:07:38--00:13:10 与 00:13:32--00:15:04 对谈内容整理。}
\end{figure}

\begin{knowledgebox}{读图：从“互联网平台收入”看到“劳动预算”}
前三层是相对看得见的收入：订阅、广告电商和 API。后两层更难估计：Agent 如果能端到端替代白领工作，就吃到工资预算；AI 如果能写 CUDA、做科学发现或寻找新材料，就创造现在没有被人类充分完成的增量价值。
\end{knowledgebox}

\subsection{看得见的收入：订阅、广告电商和 API}

广密的第一类收入是订阅。可用一个简化公式表示：
\[
R_{\text{sub}} = U \times p \times a
\]
其中，\(R_{\text{sub}}\) 表示订阅年收入；\(U\) 表示月活或可触达用户规模；\(p\) 表示付费率；\(a\) 表示单个付费用户的年均收入。节目中给出的乐观估计是：如果 ChatGPT 做到数十亿级活跃用户，且 10\% 左右用户付费，订阅收入可以进入数百亿美元量级。这一类推来自 Office、Netflix 等成熟订阅产品。

第二类收入是广告和电商。逻辑是：大流量、高频、有粘性的入口最终会变现。短视频和手机广告早期也被质疑过屏幕小、场景不适合，但后来证明用户注意力在哪里，广告和交易就会迁移到哪里。ChatGPT 如果融合搜索、推荐、购物和服务撮合，就可能从传统搜索广告和平台抽成中拿走一部分。

第三类收入是 API，也就是“模型云”。广密把它类比为 Amazon 从电商中生长出 AWS：当公司自身需要巨大计算、数据和工程系统时，这套能力也可以对外服务开发者。对 OpenAI 来说，模型能力越强、成本越低、开发者生态越大，API 收入就越像新的云业务。

\begin{importantbox}{看得见收入的共同前提}
订阅、广告电商和 API 都依赖同一件事：ChatGPT 必须保持大流量、高频使用、强品牌和足够好的模型体验。如果用户入口迁移到 Gemini、手机厂商或垂直 Agent，这三类收入的上限都会被压低。
\end{importantbox}

\subsection{看不见的收入：Agent 替代劳动}

更大的想象来自劳动预算。工具软件卖的是效率，Agent 如果真正端到端完成任务，卖的就可能是“替代一部分工资”。软件开发是最早出现这种信号的行业：AI Coding 在一年内长出百亿美元级别市场，说明知识工作中确实存在可被模型迅速吸收的任务。

但这里有两个关键门槛。第一，Agent 必须可靠到能够承担责任，不只是生成建议；第二，企业必须愿意把预算从人力、外包或传统软件迁移到 AI 系统。只有越过这两个门槛，AI 才不只是“帮员工更快”，而是成为新的劳动供给。

\begin{warningbox}{工具预算与工资预算不是一回事}
一个员工愿意每月为工具付几十美元，不等于公司愿意把一个岗位预算交给 Agent。后者要求稳定性、权限控制、审计、责任归属和可预期结果。很多 AI 公司会卡在“很好用的工具”，而不是“能吃工资预算的系统”。
\end{warningbox}

\subsection{更远的收入：AI 做人做不了的事}

广密还提到另一类看不见收入：AI 可能完成现在人类做不了或做得很慢的高价值任务，例如重写 CUDA、发现新材料、加速医学研究、寻找室温超导线索，甚至“AI 做 AI 科学家”。这类收入不能简单归入存量替代，因为它可能创造新问题、新市场和新科学路径。

这里的前提更强：模型不仅要会调用现有知识，还要在长期实验、反馈、推理、记忆和自主学习中改进。这就自然引出后文的 Online Learning。换句话说，OpenAI 的长期收入想象越大，对模型可靠性、持续学习和真实环境反馈的要求越高。

\subsection{本章小结}

OpenAI 收入账可以分成三层：订阅、广告电商和 API 是看得见的平台收入；Agent 替代白领劳动是更大的存量迁移；AI 做科学发现和高价值增量任务，则依赖更根本的模型能力突破。短期质疑来自财务账，长期乐观来自平台账和范式账。

